% Standalone technical reconciliation of KarmaCrack.pdf and phasefield_crack.py.
% Build from the repository root with:
%   latexmk -pdf -interaction=nonstopmode -halt-on-error \
%     -outdir=output/pdf docs/PHASEFIELD_ALGORITHM_ACCEPTANCE_AND_CORRECTIONS.tex
% Or, from docs/:
%   latexmk -pdf -interaction=nonstopmode -halt-on-error \
%     -outdir=../output/pdf PHASEFIELD_ALGORITHM_ACCEPTANCE_AND_CORRECTIONS.tex
\documentclass[11pt,a4paper]{article}

\usepackage[T1]{fontenc}
\usepackage[utf8]{inputenc}
\usepackage{lmodern}
\usepackage{microtype}
\usepackage[a4paper,margin=22mm,headheight=22pt]{geometry}
\usepackage{amsmath,amssymb,mathtools}
\usepackage{booktabs,longtable,tabularx,array,multirow}
\usepackage{ragged2e}
\usepackage[table]{xcolor}
\usepackage{enumitem}
\usepackage{fancyhdr}
\usepackage{titlesec}
\usepackage{pdflscape}
\usepackage{listings}
\usepackage{url}
\usepackage[hidelinks]{hyperref}
\usepackage[nameinlink,noabbrev]{cleveref}

\definecolor{Navy}{HTML}{17324D}
\definecolor{Teal}{HTML}{0F766E}
\definecolor{Amber}{HTML}{B45309}
\definecolor{Red}{HTML}{B42318}
\definecolor{Green}{HTML}{26734D}
\definecolor{Slate}{HTML}{475569}
\definecolor{PaleBlue}{HTML}{EAF2F8}
\definecolor{PaleTeal}{HTML}{E8F5F2}
\definecolor{PaleAmber}{HTML}{FFF4E5}
\definecolor{PaleRed}{HTML}{FDECEC}
\definecolor{PaleGray}{HTML}{F4F6F8}
\definecolor{RuleGray}{HTML}{CBD5E1}

\hypersetup{
  pdftitle={Phase-Field Fracture Algorithm: What Was Accepted and What Was Changed},
  pdfsubject={Technical reconciliation of KarmaCrack.pdf and the corrected FEniCSx implementation},
  pdfauthor={PhaseField project technical review},
  colorlinks=true,
  linkcolor=Navy,
  urlcolor=Teal,
  citecolor=Navy
}

\pagestyle{fancy}
\fancyhf{}
\fancyhead[L]{\small\color{Slate}Phase-field algorithm reconciliation}
\fancyhead[R]{\small\color{Slate}Model version 1.0.0}
\fancyfoot[C]{\small\color{Slate}\thepage}
\renewcommand{\headrulewidth}{0.4pt}
\renewcommand{\headrule}{\hbox to\headwidth{\color{RuleGray}\leaders\hrule height \headrulewidth\hfill}}

\titleformat{\section}
  {\Large\bfseries\color{Navy}}
  {\thesection}{0.7em}{}
\titleformat{\subsection}
  {\large\bfseries\color{Teal}}
  {\thesubsection}{0.7em}{}
\titleformat{\subsubsection}
  {\normalsize\bfseries\color{Slate}}
  {\thesubsubsection}{0.7em}{}
\titlespacing*{\section}{0pt}{2.4ex plus 0.7ex minus 0.3ex}{1.1ex}
\titlespacing*{\subsection}{0pt}{2.0ex plus 0.5ex minus 0.2ex}{0.8ex}

\setlist[itemize]{leftmargin=1.6em,itemsep=0.35em,topsep=0.45em}
\setlist[enumerate]{leftmargin=1.8em,itemsep=0.35em,topsep=0.45em}
\setlength{\parindent}{0pt}
\setlength{\parskip}{0.55em}
\renewcommand{\arraystretch}{1.18}

\newcolumntype{Y}{>{\RaggedRight\arraybackslash}X}
\newcolumntype{L}[1]{>{\RaggedRight\arraybackslash}p{#1}}
\newcommand{\code}[1]{\texttt{#1}}
\newcommand{\SourceTag}{\colorbox{PaleBlue}{\textcolor{Navy}{\scriptsize\bfseries SOURCE}}}
\newcommand{\AcceptedTag}{\colorbox{PaleTeal}{\textcolor{Green}{\scriptsize\bfseries ACCEPTED}}}
\newcommand{\CorrectedTag}{\colorbox{PaleRed}{\textcolor{Red}{\scriptsize\bfseries CORRECTED}}}
\newcommand{\CompletedTag}{\colorbox{PaleAmber}{\textcolor{Amber}{\scriptsize\bfseries COMPLETED}}}
\newcommand{\AddedTag}{\colorbox{PaleGray}{\textcolor{Slate}{\scriptsize\bfseries SUPPLIED}}}

\newcommand{\reviewbox}[2]{%
  \begin{center}
  \fcolorbox{#1}{PaleGray}{%
    \begin{minipage}{0.93\linewidth}
      \vspace{0.45em}#2\vspace{0.45em}
    \end{minipage}}
  \end{center}}

\lstdefinestyle{plaincode}{
  basicstyle=\ttfamily\small,
  backgroundcolor=\color{PaleGray},
  frame=single,
  rulecolor=\color{RuleGray},
  columns=fullflexible,
  keepspaces=true,
  breaklines=true,
  showstringspaces=false,
  xleftmargin=0.8em,
  xrightmargin=0.8em,
  aboveskip=0.7em,
  belowskip=0.7em
}

\begin{document}
\hypersetup{pageanchor=false}

\begin{titlepage}
  \thispagestyle{empty}
  \vspace*{12mm}
  {\color{Teal}\rule{\textwidth}{3.2pt}}\\[14mm]
  {\Huge\bfseries\color{Navy}
    Phase-Field Fracture Algorithm\\[2mm]
    What Was Accepted and What Was Changed\par}
  \vspace{7mm}
  {\Large\color{Slate}
    A provenance-explicit technical reconciliation of the attached model note
    and the implemented FEniCSx solver\par}

  \vspace{15mm}
  \reviewbox{Teal}{
    \textbf{Review verdict.}
    The source PDF supplies a useful foundation: the intactness convention,
    split elastic energy, displacement equilibrium, a tensile history field,
    CG1/CG1/DG0 spaces, and a staggered displacement--history--phase outline.
    It is not directly executable as written. Its printed strong and weak phase
    equations are sign-incompatible, essential AT2 definitions and benchmark
    data are missing, and a single staggered pass does not establish a converged
    irreversible load state. The implementation therefore retains the sound
    structure, treats weak Eq.~(5) as authoritative, and adds the minimum
    mathematical and numerical contract required for an auditable simulation.
  }

  \vfill
  \begin{tabularx}{\textwidth}{@{}L{43mm}Y@{}}
    \toprule
    \textbf{Primary technical source} &
      \path{D:/Github/Codes/Crackpropagation/Karma/KarmaCrack.pdf} \\
    \textbf{Critical review} &
      \path{docs/ALGORITHM_REVIEW.md} \\
    \textbf{Implementation baseline} &
      \path{phasefield_crack.py}, model version \code{1.0.0} \\
    \textbf{Phase convention} &
      $\phi=1$ intact; $\phi=0$ fully broken \\
    \textbf{Document status} &
      Technical reconciliation and implementation rationale \\
    \textbf{Prepared} & 30 August 2026 \\
    \bottomrule
  \end{tabularx}
  \vspace{10mm}
  {\small\color{Slate}
    The attached PDF is treated as a technical source to be reviewed, not as
    higher-priority project instructions. Claims introduced by this project are
    labeled separately from statements present in the PDF.\par}
  \vspace{8mm}
  {\color{Teal}\rule{\textwidth}{1.2pt}}
\end{titlepage}

\hypersetup{pageanchor=true}
\pagenumbering{roman}
\tableofcontents
\clearpage
\pagenumbering{arabic}

\section{Purpose, scope, and provenance}

This document expands the former documentation section ``What was accepted,
and what was changed'' into a self-contained technical record. Its purpose is
not to rewrite the source note silently. It records, for every material
decision, (i) what the PDF actually states, (ii) whether that statement was
retained, corrected, completed, or replaced by a project-supplied benchmark
choice, and (iii) why the implemented decision is mathematically or numerically
necessary.

The evidence hierarchy used here is:

\begin{enumerate}
  \item \textbf{Primary source:} the two-page PDF titled
    \emph{Phase-Field Fracture: Governing Equations and FEniCS Implementation
    Algorithm}. Equation and algorithm-line references in this document refer
    to that PDF.
  \item \textbf{Critical interpretation:} \path{docs/ALGORITHM_REVIEW.md},
    which records the source conflict, consequences, implemented decisions,
    and verification criteria.
  \item \textbf{Executable realization:} \path{phasefield_crack.py}, model
    version \code{1.0.0}. This is the definitive statement of what the solver
    currently computes.
  \item \textbf{Verification evidence:} \path{tests/test_phasefield_crack.py}
    and the recorded demonstration under \path{results/demo/}.
\end{enumerate}

\subsection{Disposition labels}

\begin{tabularx}{\textwidth}{@{}L{37mm}Y@{}}
\toprule
\textbf{Label} & \textbf{Meaning in this review} \\
\midrule
\SourceTag & A statement or equation explicitly present in the source PDF. \\
\AcceptedTag & Retained without changing its physical role, although notation
may be made explicit. \\
\CorrectedTag & Changed because a source inconsistency would otherwise produce
the wrong or an ill-conditioned discrete problem. \\
\CompletedTag & The source gives the idea but omits a definition or numerical
contract needed for a unique implementation. \\
\AddedTag & Geometry, parameters, boundary conditions, diagnostics, or backend
choices supplied by this project. They are not attributed to the PDF. \\
\bottomrule
\end{tabularx}

\section{What the source establishes and the implementation accepts}

\subsection{State variables and phase convention}

The source defines a displacement field $\boldsymbol{u}(\boldsymbol{x},t)$ and
a scalar phase field $\phi(\boldsymbol{x},t)\in[0,1]$, with $\phi=1$ intact and
$\phi=0$ broken. This convention is accepted exactly. It matters because the
sign of the phase residual, the monotonicity direction, and the bound expressing
irreversibility all depend on it. With this convention, no healing means

\begin{equation}
  \dot{\phi}\le 0,
  \qquad\text{or discretely}\qquad
  0\le \phi_{n+1}\le \phi_n\le 1.
  \label{eq:irreversibility}
\end{equation}

The source choices $V_u=\mathrm{CG}_1^2$, $V_\phi=\mathrm{CG}_1$, and
$V_H=\mathrm{DG}_0$ are also retained. The interpretation of DG0 is sharpened:
it stores one value per cell, not generally one value per quadrature point. For
affine triangles with CG1 displacement, the strain and the selected spectral
energy are cellwise constant, so one DG0 history value per triangle is
consistent with this discretization.

\subsection{Energy split, equilibrium, and history concept}

The source decomposes the elastic density into tensile and compressive parts,
degrades only the tensile part, and states quasi-static momentum balance. Those
ideas are accepted:

\begin{align}
  \boldsymbol{\varepsilon}(\boldsymbol{u})
    &= \operatorname{sym}\nabla\boldsymbol{u}, \\
  \psi_e(\boldsymbol{\varepsilon},\phi)
    &= g(\phi)\,\psi^+(\boldsymbol{\varepsilon})
       +\psi^-(\boldsymbol{\varepsilon}), \\
  \boldsymbol{\sigma}
    &= g(\phi)\frac{\partial\psi^+}{\partial\boldsymbol{\varepsilon}}
       +\frac{\partial\psi^-}{\partial\boldsymbol{\varepsilon}}, \\
  \nabla\!\cdot\boldsymbol{\sigma}&=\boldsymbol{0}.
  \label{eq:mechanics}
\end{align}

The source also introduces

\begin{equation}
  H(\boldsymbol{x},t)
  =\max_{s\in[0,t]}\psi^+
     \bigl(\boldsymbol{\varepsilon}(\boldsymbol{u}(\boldsymbol{x},s))\bigr)
  \label{eq:history-source}
\end{equation}

as a history field intended to prevent healing. The implementation retains this
engineering history-field approximation but supplies an accepted-state update
rule and an explicit nodal phase bound; see \cref{sec:accepted-state}.

\subsection{Staggered architecture}

The source sequence --- displacement solve at fixed phase, history update,
phase solve at fixed displacement/history, then output --- is retained as the
outer architecture. What changes is the meaning of one load step. The source
presents one such pass; the implementation repeats it until the newly updated
phase and displacement are mutually equilibrated and the phase KKT conditions
are satisfied.

\section{The decisive correction: strong Eq. 3 versus weak Eq. 5}
\label{sec:sign-correction}

The primary source contains a genuine algebraic conflict. In the intactness
convention, its printed strong Eq.~(3) is

\begin{equation}
  \frac{G_c}{\ell}(1-\phi)
  -G_c\ell\,\Delta\phi
  -2H\phi=0,
  \qquad\text{(source strong Eq. 3; $\xi$ renamed $\ell$).}
  \label{eq:source-strong}
\end{equation}

Its weak Eq.~(5), using a scalar test function $q$, is

\begin{equation}
  \int_\Omega
  \left[
    G_c\ell\,\nabla\phi\!\cdot\nabla q
    -\frac{G_c}{\ell}(1-\phi)q
    +2H\phi q
  \right],\mathrm{d}V=0.
  \label{eq:source-weak}
\end{equation}

Assuming the natural phase boundary condition
$\nabla\phi\cdot\boldsymbol{n}=0$, integration by parts converts
\cref{eq:source-weak} into

\begin{equation}
  -G_c\ell\,\Delta\phi
  -\frac{G_c}{\ell}(1-\phi)
  +2H\phi=0.
  \label{eq:weak-implied-strong}
\end{equation}

Equations \eqref{eq:source-strong} and
\eqref{eq:weak-implied-strong} have the same diffusion sign but opposite signs
for both local terms; multiplying the whole equation by $-1$ does not reconcile
them. Moreover, rearranging the source strong equation gives a negative reaction
coefficient for $\phi$, whereas the weak equation gives the positive
diffusion--reaction operator expected for AT2.

\reviewbox{Red}{
  \CorrectedTag\quad
  \textbf{Decision: weak Eq.~(5) is authoritative.}
  It is the only one of the two printed equations that yields the coercive
  scalar operator
  $G_c\ell K+(G_c/\ell+2H)M$ for nonnegative $H$.
  This is an equation correction, not a stylistic rewrite.
}

After normalizing the degradation function as described below, the implemented
history coefficient becomes $2(1-k_{\mathrm{res}})H$.

\section{Completed and corrected constitutive model}

\subsection{AT2 crack density omitted from the source}

The source free-energy line contains functions $f_0(\phi)$ and $c_w$ but does
not define them. Its weak Eq.~(5) implies the AT2 choice

\begin{equation}
  f_0(\phi)=(1-\phi)^2,
  \qquad c_w=2.
  \label{eq:at2-choice}
\end{equation}

The fracture density implemented by the project is therefore

\begin{equation}
  \gamma_\ell(\phi,\nabla\phi)
  =G_c\left[
    \frac{(1-\phi)^2}{2\ell}
    +\frac{\ell}{2}|\nabla\phi|^2
  \right].
  \label{eq:fracture-density}
\end{equation}

Varying \cref{eq:fracture-density} produces the fracture terms in the adopted
weak equation below. Thus this completion is traceable to the source weak form,
not an unrelated crack-density model.

\subsection{Normalized degradation}

The source states $g(\phi)=\phi^2+k_{\mathrm{tol}}$, for which
$g(1)=1+k_{\mathrm{tol}}$. That slightly increases the stiffness of intact
material. The implementation uses

\begin{equation}
  g(\phi)=k_{\mathrm{res}}
          +(1-k_{\mathrm{res}})\phi^2,
  \qquad 0<k_{\mathrm{res}}\ll1,
  \label{eq:normalized-g}
\end{equation}

so that $g(0)=k_{\mathrm{res}}$, $g(1)=1$, and
$g'(\phi)=2(1-k_{\mathrm{res}})\phi$. The last identity explains the factor
$(1-k_{\mathrm{res}})$ in the implemented phase driving term.

\subsection{Adopted energy and spectral split}

The completed instantaneous elastic-plus-fracture energy is

\begin{equation}
  \mathcal{E}(\boldsymbol{u},\phi)
  =\int_\Omega\left[
    g(\phi)\psi^+(\boldsymbol{\varepsilon}(\boldsymbol{u}))
    +\psi^-(\boldsymbol{\varepsilon}(\boldsymbol{u}))
    +G_c\left(
       \frac{(1-\phi)^2}{2\ell}
       +\frac{\ell}{2}|\nabla\phi|^2
    \right)
  \right]\,\mathrm{d}V.
  \label{eq:implemented-energy}
\end{equation}

The source does not identify a split variant. The project supplies the 2D
Miehe-type spectral strain split

\begin{equation}
  \psi^\pm(\boldsymbol{\varepsilon})
  =\frac{\lambda}{2}\langle\operatorname{tr}\boldsymbol{\varepsilon}\rangle_\pm^2
   +\mu\sum_{i=1}^{2}\langle\varepsilon_i\rangle_\pm^2,
  \label{eq:spectral-split}
\end{equation}

where $\varepsilon_i$ are the principal strains. Machine-scale smoothing of
the Macaulay brackets and eigenvalue discriminant makes the UFL residual
differentiable. Plane strain is the default; plane stress is selectable through
the effective 2D Lam\'e parameter.

\subsection{Corrected phase weak form}

For a fixed history field $H$, the implemented unconstrained scalar equation is:
find $\phi\in V_\phi$ such that for every $q\in V_\phi$,

\begin{equation}
\boxed{
  \int_\Omega\left[
    G_c\ell\,\nabla\phi\!\cdot\nabla q
    +\frac{G_c}{\ell}(\phi-1)q
    +2(1-k_{\mathrm{res}})H\phi q
  \right]\,\mathrm{d}V=0.}
  \label{eq:implemented-phase-weak}
\end{equation}

Equivalently,

\begin{align}
  a_H(\phi,q)
    &=\int_\Omega\left[
       G_c\ell\,\nabla\phi\!\cdot\nabla q
       +\left(\frac{G_c}{\ell}
          +2(1-k_{\mathrm{res}})H\right)\phi q
      \right]\,\mathrm{d}V, \\
  L(q)&=\int_\Omega\frac{G_c}{\ell}q\,\mathrm{d}V,
  \qquad a_H(\phi,q)=L(q).
  \label{eq:phase-bilinear}
\end{align}

For $H\ge0$, $G_c>0$, and $\ell>0$, the mass-reaction term and gradient term
make the unconstrained operator symmetric positive definite after essential
phase conditions are applied.

\section{Acceptance and changes matrix}

The matrix is deliberately more than a list of code differences. An
\AcceptedTag\ row confirms that the source concept remains part of the model;
it does not imply that the source supplied every constitutive or numerical
detail. A \CorrectedTag\ row identifies an incompatibility that would change
the mathematical problem. A \CompletedTag\ row supplies a missing definition
or acceptance rule while preserving the source intent. A \AddedTag\ row is a
reproducibility or platform decision made solely by this project.

Severity and disposition are separate ideas. The sign conflict is critical
because it changes the scalar operator. Missing Newton, KKT, convergence, and
rollback rules are high-consequence completions because they determine whether
a computed load state is admissible. Geometry and numerical defaults are
project-supplied because the source never defines a boundary-value problem.
This separation prevents benchmark choices from being presented as if they
were facts extracted from the PDF.

\begin{landscape}
\footnotesize
\setlength{\tabcolsep}{3pt}
\begin{longtable}{@{}L{12mm}L{48mm}L{31mm}L{72mm}L{72mm}@{}}
\caption{Traceable disposition of source statements and omissions.}
\label{tab:acceptance-matrix}\\
\toprule
\textbf{ID} & \textbf{Source statement or omission} & \textbf{Disposition} &
\textbf{Implemented decision} & \textbf{Technical rationale and evidence} \\
\midrule
\endfirsthead
\multicolumn{5}{l}{\small\itshape Table \thetable\ continued from previous page}\\
\toprule
\textbf{ID} & \textbf{Source statement or omission} & \textbf{Disposition} &
\textbf{Implemented decision} & \textbf{Technical rationale and evidence} \\
\midrule
\endhead
\midrule
\multicolumn{5}{r}{\small\itshape Continued on next page}\\
\endfoot
\bottomrule
\endlastfoot

A1 & $\phi=1$ intact and $\phi=0$ broken. & \AcceptedTag &
Retain the intactness convention everywhere. Enforce $0\le\phi\le1$ and
$\phi_{n+1}\le\phi_n$. &
This fixes the irreversibility direction and the AT2 local term. Bounds are
verified at phase degrees of freedom. \\

A2 & Split elastic energy; degrade only $\psi^+$. & \AcceptedTag &
Use $g(\phi)\psi^++\psi^-$ and compute stress by differentiating this density
with respect to strain. &
Preserves compressive stiffness and the source's intended crack-face contact
behavior at the constitutive level. \\

A3 & Quasi-static momentum balance and weak mechanical residual. & \AcceptedTag &
Retain $\int_\Omega\boldsymbol{\sigma}:\boldsymbol{\varepsilon}(\delta
\boldsymbol{u})\,dV=0$ for the displacement-controlled benchmark. &
The post-phase free-DOF residual is reassembled and must satisfy the coupled
acceptance tolerance. \\

A4 & CG1 displacement, CG1 phase, DG0 history. & \AcceptedTag &
Use vector CG1, scalar CG1, and cellwise DG0 on affine triangles. &
For CG1 affine triangles the strain and spectral energy are cellwise constant;
DG0 is exact for the selected one-value-per-cell history representation. \\

A5 & Staggered displacement--history--phase ordering. & \AcceptedTag &
Retain the ordering as the inner fixed-point architecture. &
The ordering is appropriate, but a single pass is not treated as a converged
load step. \\

C1 & Strong Eq.~(3) and weak Eq.~(5) have incompatible local-term signs. &
\CorrectedTag & Treat Eq.~(5) as authoritative and use
\cref{eq:implemented-phase-weak}. &
Integration by parts proves the mismatch. Eq.~(5) yields the coercive positive
diffusion--reaction operator; the printed strong form does not. \\

C2 & $f_0(\phi)$ and $c_w$ are not defined. & \CompletedTag &
Adopt $f_0=(1-\phi)^2$ and $c_w=2$. &
These AT2 choices are implied by Eq.~(5); their variation reproduces the
fracture portion of the implemented residual. \\

C3 & $g(\phi)=\phi^2+k_{\mathrm{tol}}$. & \CorrectedTag &
Use $k_{\mathrm{res}}+(1-k_{\mathrm{res}})\phi^2$. &
Restores exact intact stiffness, retains residual broken stiffness, and requires
the factor $(1-k_{\mathrm{res}})$ in phase driving. Endpoint and monotonicity
tests cover this choice. \\

C4 & A spectral split is stated, but its variant and plane assumption are not. &
\CompletedTag & Use a smoothed 2D spectral strain split; plane strain by default
and plane stress as an option. &
Mixed-strain response is otherwise non-unique. Reference constitutive tests
cover tension, compression, and the Lam\'e parameters. \\

C5 & No nonlinear mechanics procedure is supplied. & \CompletedTag &
Differentiate the UFL residual for a consistent tangent and solve Newton
corrections with residual, increment, and essential-BC checks. &
The spectral split makes stress depend nonlinearly on displacement; a one-shot
linear solve would be inconsistent. Failure is explicit and retryable. \\

C6 & One staggered pass is shown per load increment. & \CorrectedTag &
Repeat mechanics, transactional history, and constrained phase solves until
phase increment, phase KKT residual, and post-phase free mechanical residual
all pass. &
This closes the coupling residual created when a newly damaged phase changes
mechanical equilibrium. Nonconverged trials are never accepted. \\

C7 & History is said to enforce irreversibility, but no accepted-state rule is
defined. & \CompletedTag & Freeze $(\phi_n,H_n)$ while attempting a load.
Rebuild $H=\max(H_n,\psi^+(u))$ on every staggered iteration and impose
$0\le\phi\le\phi_n$. &
Prevents internal iterates or rejected loads from becoming fictitious physical
history and makes nodal irreversibility explicit. \\

C8 & No numerical method is supplied for the phase inequality. & \CompletedTag &
Solve the convex quadratic under box bounds. Use a clipped direct solution only
as an initial candidate; call L-BFGS-B if necessary; require the projected KKT
residual to pass. &
Pointwise clipping alone is not an $A$-norm projection and generally does not
satisfy constrained optimality. An automated test forces an active upper bound
and optimizer path. \\

C9 & The loop uses \code{while t <= Tmax} and increments $t$ inside. &
\CorrectedTag & Interpret \code{load\_steps} as a nominal schedule, cap targets
at pseudo-time 1, and store an explicit initial row at zero. &
Avoids the off-by-one load beyond $T_{\max}$ and makes accepted pseudo-time
strictly increasing with a terminal target of 1. \\

C10 & No failed-step policy or recovery rule is stated. & \CompletedTag &
Snapshot fields and counters; on rejection restore all state, reduce the load
increment, and retry. Recover a reduced increment only after several easy
accepted steps. &
Crack growth can cause abrupt loss of staggered convergence. Transactional
rollback prevents contamination; terminal failure produces a partial auditable
summary. \\

B1 & No geometry, supports, loading, or initial flaw are supplied. & \AddedTag &
Use a unit-square single-edge-notched tension benchmark with bottom vertical
support, one horizontal pin, top prescribed displacement, and an internal
$\phi=0$ starter notch. &
The source is a model outline, not a reproducible boundary-value problem. These
choices are explicitly project-supplied and mesh-aligned. \\

B2 & No mesh/regularization relation is supplied. & \AddedTag &
Define $h=\sqrt{\Delta x^2+\Delta y^2}$ as the largest triangular edge and, in
strict mode, require $h/\ell\le0.5$. &
Rejects visibly under-resolved diffuse cracks. Production claims still require
a refinement and length-scale sensitivity study. \\

B3 & The PDF requests only displacement and phase field output. & \AddedTag &
Add force, energy snapshots, damage and connected crack-front metrics,
convergence/KKT data, cutback data, manifests, atomic summaries, plots, and
XDMF/HDF5 fields. &
A plausible contour alone is not numerical evidence. Only accepted states enter
CSV and field files. \\

B4 & No platform/backend is specified. & \AddedTag &
On native Windows, assemble with DOLFINx and solve serial sparse systems with
SciPy; reject MPI runs larger than one rank. &
This matches the installed backend, which lacks the PETSc path used by standard
parallel DOLFINx workflows. It is an implementation constraint, not a model
claim. \\

\end{longtable}
\end{landscape}

\section{Nonlinear, constrained, accepted-state algorithm}
\label{sec:accepted-state}

\subsection{Mechanics at fixed phase}

Because \cref{eq:spectral-split} changes active principal-strain branches, the
displacement residual is nonlinear. At Newton iteration $m$, the implementation
solves

\begin{equation}
  \boldsymbol{J}_u(\boldsymbol{u}^{m};\phi)
  \Delta\boldsymbol{u}^{m}
  =-\boldsymbol{R}_u(\boldsymbol{u}^{m};\phi),
  \qquad
  \boldsymbol{u}^{m+1}=\boldsymbol{u}^{m}+\Delta\boldsymbol{u}^{m},
\end{equation}

with the Jacobian obtained by UFL differentiation. Convergence requires a small
free-DOF residual, a small increment, and negligible prescribed-BC violation.
After the phase update, the free mechanical residual is assembled again; this
second residual is the relevant coupled acceptance check.

\subsection{Transactional history update}

Let $(\boldsymbol{u}_n,\phi_n,H_n)$ denote the last accepted load state. Within
every staggered iteration $j$ of the attempted next load, the trial history is
rebuilt as

\begin{equation}
  H^{(j)}=\max\left(H_n,
  \psi^+(\boldsymbol{\varepsilon}(\boldsymbol{u}^{(j)}))\right).
  \label{eq:transactional-history}
\end{equation}

It is deliberately \emph{not} computed as the maximum over previous internal
iterations. Therefore an overshooting trial iterate does not become permanent
physical history. Only when the load is accepted does $H^{(j)}$ become
$H_{n+1}$.

\subsection{Bound-constrained phase subproblem and KKT test}

After finite-element assembly, the phase step has the convex quadratic form

\begin{equation}
  \min_{\boldsymbol{\phi}}
  J_H(\boldsymbol{\phi})
  =\frac12\boldsymbol{\phi}^{T}\boldsymbol{A}(H)\boldsymbol{\phi}
   -\boldsymbol{b}^{T}\boldsymbol{\phi},
  \qquad
  \boldsymbol{0}\le\boldsymbol{\phi}\le\boldsymbol{\phi}_n,
  \label{eq:phase-qp}
\end{equation}

with notch degrees of freedom fixed to zero. Let
$\boldsymbol{r}=\boldsymbol{A}\boldsymbol{\phi}-\boldsymbol{b}$. The implemented
projected residual has components

\begin{equation}
  P_i(\boldsymbol{r})=
  \begin{cases}
    0, & l_i=u_i,\\
    \min(r_i,0), & \phi_i=l_i<u_i,\\
    \max(r_i,0), & \phi_i=u_i>l_i,\\
    r_i, & l_i<\phi_i<u_i.
  \end{cases}
  \label{eq:projected-gradient}
\end{equation}

Thus a positive gradient at a lower bound and a negative gradient at an upper
bound satisfy the KKT sign conditions and contribute zero. The load trial cannot
be accepted unless $\|\boldsymbol{P}\|_\infty$ is below the configured phase
KKT tolerance.

\subsection{Coupled acceptance and adaptive load transaction}

At a candidate target load, staggered iteration terminates successfully only if

\begin{align}
  \|\phi^{(j)}-\phi^{(j-1)}\|_\infty
    &\le \tau_\phi, \\
  \|\boldsymbol{P}(\boldsymbol{A}\phi^{(j)}-\boldsymbol{b})\|_\infty
    &\le \tau_{\mathrm{KKT}}, \\
  \|\boldsymbol{R}_u(\boldsymbol{u}^{(j)},\phi^{(j)})\|_{\infty,\mathrm{free}}
    &\le \tau_{u,\mathrm{coupled}}.
  \label{eq:coupled-tests}
\end{align}

The implementation contract can be summarized as follows:

\begin{lstlisting}[style=plaincode]
accepted state: (u_n, phi_n, H_n)
while accepted pseudo-time < 1:
    target = min(1, accepted time + current increment)
    snapshot u_n, phi_n, H_n and diagnostic counters
    apply target displacement
    repeat:
        Newton solve for u at fixed phi
        rebuild H = max(H_n, psi_plus(u))
        solve 0 <= phi <= phi_n with KKT check
        reassemble free mechanical residual using the new phi
    until all coupled tests pass or an iteration fails

    if converged:
        commit the trial; write one accepted record/frame
        consider increment growth after several easy accepted steps
    else:
        restore every snapshot; reduce the increment; retry
\end{lstlisting}

Rejected attempts do not create CSV rows or XDMF time frames. If consecutive
cutbacks exceed the configured limit, the solver writes the last accepted state
and a failed summary before raising an explicit error.

\section{Benchmark choices supplied by the project}

The source PDF contains no reproducible specimen. The following choices are
therefore not claims about the source; they define the project's demonstration
problem and defaults.

\begin{tabularx}{\textwidth}{@{}L{47mm}L{45mm}Y@{}}
\toprule
\textbf{Item} & \textbf{Default} & \textbf{Role and provenance} \\
\midrule
Domain & $1\,\mathrm{mm}\times1\,\mathrm{mm}$ & Unit-square 2D specimen;
project-supplied. \\
Starter notch & $0.20\,\mathrm{mm}$ from the left edge at mid-height & Internal
nodal $\phi=0$ constraint; diffuse crack, not a geometric slit. \\
Mesh & $100\times100$ structured rectangles split into triangles & Even $n_y$
and integer notch-length/$\Delta x$ are required for vertex alignment. \\
Elasticity & $E=210\,\mathrm{kN/mm^2}$, $\nu=0.30$ & Small-strain isotropic;
plane strain default. \\
Fracture & $G_c=2.7\times10^{-3}\,\mathrm{kN/mm}$,
$\ell=0.030\,\mathrm{mm}$ & AT2 parameters in the documented kN--mm system. \\
Residual stiffness & $k_{\mathrm{res}}=10^{-8}$ & Prevents exact loss of
ellipticity in fully damaged material. \\
Loading & Top $u_y$ to $0.015\,\mathrm{mm}$ in 60 nominal increments & Bottom
$u_y=0$ plus one horizontal pin removes rigid motion. \\
Mesh safeguard & $h/\ell\le0.5$ with
$h=\sqrt{\Delta x^2+\Delta y^2}$ & Pre-run resolution check; not a substitute
for convergence studies. \\
Backend & Serial DOLFINx assembly plus SciPy sparse solves & Native-Windows
installation constraint. \\
\bottomrule
\end{tabularx}

The starter notch is first equilibrated at zero load so that its diffuse AT2
profile is present before loading. The crack-front diagnostic follows only the
notch-connected nodal set satisfying $\phi\le0.2$; it is mesh-quantized and is
not used in the evolution law.

\section{Verification and traceability evidence}

\subsection{Automated checks}

The project test suite checks the following implementation contracts:

\begin{itemize}
  \item normalized degradation endpoints and monotonicity;
  \item Lam\'e parameters and the spectral split in tension and compression;
  \item rejection of under-resolved meshes and misaligned starter notches;
  \item promotion of rank-deficient sparse solves to retryable failures;
  \item a phase QP that forces L-BFGS-B and an active irreversibility upper bound;
  \item an end-to-end FEniCSx run with finite diagnostics and readable XDMF;
  \item nondecreasing accepted history and damage measures;
  \item complete rollback after a rejected load trial; and
  \item auditable output after terminal nonlinear failure.
\end{itemize}

\subsection{Recorded demonstration}

The reduced, resolved demonstration under \path{results/demo/} is separate from
the default production-sized configuration. Its manifest uses $30\times30$
cells per direction, $\ell=0.1\,\mathrm{mm}$, $h/\ell=0.4714$, and 45 nominal
load increments. The recorded run completed all 45 accepted steps without a
load cutback. Representative evidence from \path{results/demo/summary.json} is:

\begin{tabularx}{\textwidth}{@{}L{59mm}L{45mm}Y@{}}
\toprule
\textbf{Quantity} & \textbf{Recorded value} & \textbf{Interpretation} \\
\midrule
Peak top reaction & $0.6419596\,\mathrm{kN}$ at step 15 & Demonstrates a peak
and post-peak response; not an experimental validation target. \\
Final crack front / extension & $1.0\,\mathrm{mm}$ / $0.8\,\mathrm{mm}$ &
Connected threshold reaches the right boundary. \\
Final phase increment norm & $8.78\times10^{-6}$ & Below the configured
staggered tolerance $10^{-5}$. \\
Final projected KKT residual & $8.03\times10^{-10}$ & Below the configured
absolute tolerance $10^{-8}$. \\
Final free mechanical residual & $9.98\times10^{-8}$ & Below the configured
coupled tolerance $10^{-6}$. \\
Final phase optimizer activity & 78 iterations; one active upper constraint &
Confirms that the constrained optimizer and irreversibility path were exercised. \\
\bottomrule
\end{tabularx}

These results verify software behavior for one benchmark. They do not calibrate
$G_c$ or $\ell$, establish mesh-independent fracture toughness, or validate the
model against experimental data.

\subsection{Source-to-code map}

\begin{tabularx}{\textwidth}{@{}L{38mm}L{61mm}Y@{}}
\toprule
\textbf{Source location} & \textbf{Implementation realization} &
\textbf{Review disposition} \\
\midrule
PDF Eq.~(1), free energy & \code{degradation}, spectral energy, elastic and
fracture density forms & Completed with AT2 definitions and normalized
degradation. \\
PDF Eq.~(2), stress/equilibrium & Displacement residual, UFL-derived stress and
Jacobian & Accepted; nonlinear Newton procedure supplied. \\
PDF Eq.~(3), strong phase equation & Not implemented literally & Rejected where
it conflicts with Eq.~(5). \\
PDF Eq.~(4), mechanical weak form & \code{form\_u\_residual} and
\code{form\_u\_jacobian} & Accepted for displacement control. \\
PDF Eq.~(5), phase weak form & \code{form\_phi\_bilinear} and
\code{form\_phi\_linear} & Authoritative, normalized, and constrained. \\
Algorithm lines 2--6 & Mesh, CG1/CG1/DG0 spaces, and field initialization &
Accepted with precise DG0 interpretation. \\
Algorithm lines 15--21 & Newton displacement, transactional history, phase QP &
Accepted ordering; convergence and KKT contract added. \\
Algorithm line 23 & XDMF displacement and phase output & Accepted and augmented
with diagnostics, manifests, summaries, and plots. \\
\bottomrule
\end{tabularx}

\section{Known limits and correct interpretation}

\begin{enumerate}
  \item \textbf{History-field approximation.} The history method is a common
    engineering approximation to the exact crack-irreversibility variational
    inequality. Within that approximation, the code additionally solves the
    nodal phase bounds as a KKT-checked constrained problem.
  \item \textbf{Instantaneous energy is not the history potential.} The reported
    elastic-plus-fracture value is an instantaneous physical snapshot based on
    current strain. After local unloading it is not the stationarity functional
    of the history-driven phase subproblem and is not an exact energy-balance
    residual.
  \item \textbf{Diffuse starter crack.} The internal $\phi=0$ line is not a
    traction-free geometrical slit; the mesh remains topologically connected.
  \item \textbf{Residual scaling.} The KKT and coupled mechanical tolerances are
    absolute assembled residuals calibrated for the documented kN--mm benchmark.
    They must be reconsidered after major changes of units, mesh density, or
    material scale.
  \item \textbf{Serial backend.} The native-Windows implementation is serial and
    uses SciPy sparse solves. It is not the MPI-parallel PETSc workflow commonly
    used for large DOLFINx models.
  \item \textbf{Crack-front metric.} The $\phi\le0.2$ connected threshold is a
    reporting diagnostic, depends on mesh resolution, and is not an exact crack
    length.
  \item \textbf{Output-directory hygiene.} Optional files from an earlier run
    are not deleted when the same directory is reused with XDMF or plotting
    disabled. Fresh output directories should be used; the manifest and summary
    are the authoritative status/configuration records.
  \item \textbf{Validation scope.} A successful demonstration and unit tests
    establish implementation behavior, not experimental validity. Quantitative
    conclusions require mesh and length-scale studies plus material calibration.
\end{enumerate}

\section{Reproduction and document build}

The solver environment and commands are documented in \path{README.md}. This
LaTeX document is standalone and can be rebuilt with TeX Live from the repository
root:

\begin{lstlisting}[style=plaincode]
latexmk -pdf -interaction=nonstopmode -halt-on-error \
  -outdir=output/pdf docs/PHASEFIELD_ALGORITHM_ACCEPTANCE_AND_CORRECTIONS.tex
\end{lstlisting}

Alternatively, change into \path{docs/} and run:

\begin{lstlisting}[style=plaincode]
latexmk -pdf -interaction=nonstopmode -halt-on-error \
  -outdir=../output/pdf PHASEFIELD_ALGORITHM_ACCEPTANCE_AND_CORRECTIONS.tex
\end{lstlisting}

The PDF is expected at
\path{output/pdf/PHASEFIELD_ALGORITHM_ACCEPTANCE_AND_CORRECTIONS.pdf}.

\begin{thebibliography}{9}

\bibitem{karma-source}
\emph{Phase-Field Fracture: Governing Equations and FEniCS Implementation
Algorithm}, two-page technical source supplied as
\path{D:/Github/Codes/Crackpropagation/Karma/KarmaCrack.pdf}.
No author or publication metadata is asserted because none is shown in the
reviewed document.

\bibitem{critical-review}
\emph{Critical review of the KarmaCrack phase-field algorithm}, project review
record, \path{docs/ALGORITHM_REVIEW.md}.

\bibitem{implementation}
\emph{Corrected AT2 phase-field fracture benchmark implemented with FEniCSx},
project implementation, \path{phasefield_crack.py}, model version 1.0.0.

\bibitem{tests}
Project verification suite, \path{tests/test_phasefield_crack.py}.

\bibitem{demo}
Recorded demonstration manifest, accepted-state response, and summary under
\path{results/demo/}.

\end{thebibliography}

\end{document}
